\begin{lemma}
In a $3$-uniform $3$-simple local extremal setup at scale $D$, with
$\mathcal F$, $f$, $X_s$, and $X_b$ as in
\noderef{ThreeUniformThreeSimpleLocalSetup}, if
$P(f)=(1+\delta)D^2-9D$ with $0\le\delta\le0.031$, then the high-degree part
satisfies
\[
|X_b|=6.
\]
\end{lemma}
\begin{proof}
The hypotheses of the lemma give
\[
P(f)=(1+\delta)D^2-9D,\qquad 0\le\delta\le0.031.
\]
The local setup \noderef{ThreeUniformThreeSimpleLocalSetup} records the exact
pair-count identity
\[
P(f)=3D^2-6D+3-w(X)+Y
\]
and the nonnegativity $Y\ge0$.  Substituting the displayed hypothesis for
$P(f)$ and solving for $w(X)$ gives
\[
w(X)=3D^2-6D+3+Y-P(f)
     =(2-\delta)D^2+3D+3+Y .
\]
Since $D$ is in the sufficiently large range of the local setup, $D$ is
positive, and therefore $3D+3+Y\ge0$.  Hence
\[
w(X)\ge (2-\delta)D^2. \tag{1}
\]

The same setup states that $X_b=X\setminus X_s$, that $X_s$ is the maximal
initial segment of the vertices of lowest total degree into $f$ subject to
\[
\sum_{v\in f}\sum_{x\in X_s}\deg_{\mathcal F}(v,x)\le D/2.
\]
It also records the row and column conventions for the table
$a_{v,x}=\deg_{\mathcal F}(v,x)$: every column sum is at most $D$, each row
sum is $2(D-1)$, and hence
\[
\sum_{v\in f}\sum_{x\in X}\deg_{\mathcal F}(v,x)=6(D-1).
\]
These are precisely the $k=3$ saturated edge table properties stated in
\noderef{ThreeUniformSaturatedEdgeStructure}, incorporated into the local
setup.  Therefore the total degree into $X_b$ is
\[
\sum_{v\in f}\sum_{x\in X_b}\deg_{\mathcal F}(v,x)
=6(D-1)-\sum_{v\in f}\sum_{x\in X_s}\deg_{\mathcal F}(v,x)
\ge 6(D-1)-D/2.
\]
For sufficiently large $D$ this last quantity is greater than $5D$.  If
$|X_b|\le5$, the column bound $D$ would make the same total at most $5D$.
Thus $|X_b|\ge6$.

Suppose, for contradiction, that $|X_b|\ge7$.  Choose $x'\in X_b$ minimizing
\[
m=\sum_{v\in f}\deg_{\mathcal F}(v,x').
\]
Since $D^2\ge0$ and $\delta\le0.031$, the present pair-count equality gives
$P(f)=(1+\delta)D^2-9D\le1.031D^2-9D$.
By \noderef{ThreeUniformThreeSimpleSmallXs}, applied with this cutoff,
the present equality $P(f)=(1+\delta)D^2-9D$, and $\delta\ge0$,
\[
\sum_{v\in f}\sum_{x\in X_s}\deg_{\mathcal F}(v,x)
\le \frac85\delta D\le0.0496D,
\]
where the second inequality is $\delta\le0.031$.  The maximality condition on
$X_s$ in \noderef{ThreeUniformThreeSimpleLocalSetup} now gives the lower
bound on any vertex of $X_b$: if some $x\in X_b$ had total degree into $f$
less than
\[
D/2-\sum_{v\in f}\sum_{y\in X_s}\deg_{\mathcal F}(v,y),
\]
then adding $x$ to the initial low-degree segment would still keep the total
at most $D/2$, contradicting maximality.  Hence, in particular,
\[
m\ge D/2-\sum_{v\in f}\sum_{y\in X_s}\deg_{\mathcal F}(v,y)
\ge D/2-0.0496D>0.45D. \tag{2}
\]
Since $x'$ is the minimum among at least seven vertices of $X_b$ and the
total over $X_b$ is at most the total over all of $X$, namely $6(D-1)$, we
also have
\[
m\le \frac{6(D-1)}{|X_b|}\le \frac67D.
\tag{3}
\]

Apply \noderef{TablePairFillingBound} to the table on $X\setminus\{x'\}$ with
three rows, column bound $D$, and total-sum bound $6(D-1)-m$.  By (2) the
parameter $m$ is positive, and by (3), in the sufficiently large range,
$6(D-1)-m=5D+(D-6-m)$ has quotient $q=5$ and a nonnegative remainder
$D-6-m\le D-m<D$.  The bound therefore gives
\[
w(X\setminus\{x'\})\le
 \frac13\bigl(5D^2+(D-6-m)^2\bigr)
\le \frac13\bigl(5D^2+(D-m)^2\bigr)
=2D^2-\frac23mD+\frac13m^2 .
\]
Applying the same bound to the single column $x'$ with total sum and column
bound both equal to $m$ gives $w(x')\le m^2/3$.  Since the local setup defines
$w(S)$ as the sum of the column contributions over $S$, the two disjoint
parts $X\setminus\{x'\}$ and $\{x'\}$ add to $X$, so
\[
w(X)=w(X\setminus\{x'\})+w(x')
\le2D^2-\frac23m(D-m).
\]
The interval (2)--(3) implies $D-m\ge D/7$, and therefore
\[
w(X)\le 2D^2-\frac23(0.45D)(D/7)
\le(2-0.04)D^2.
\]
This contradicts (1), because $\delta\le0.031<0.04$ gives
$(2-\delta)D^2>(2-0.04)D^2$ for the positive $D$ in the local setup.
Therefore $|X_b|\le6$, and together with $|X_b|\ge6$ this proves
$|X_b|=6$.
\end{proof}
